%% kapa: exam=1 round=27 year=2016 period=1 form=A subject=real_estate subject_ko=부동산학원론 seq=13 number=93
%% source: Q-Net 감정평가사 기출문제 (cst00309) artlSeq=5169402 file=extracted/1008/5169402/제27회 감정평가사 1차 문제/제27회 감정평가사 1교시문제(A형).pdf
%% conversion: kordoc 4.24.51 -> tools/gen1.py (auto); review=proofread-sonnet; flags=none
\begin{question}{93}
주택저당대출방식에 관한 설명으로 옳지 않은 것은?
\begin{choices}
  \item 원금균등분할상환방식은 대출기간 동안 매기 원금을 균등하게 분할 상환하고 이자는 점차적으로 감소하는 방식이다.
  \item 원리금균등분할상환방식의 원리금은 대출금에 감채기금계수를 곱하여 산출한다.
  \item 만기일시상환방식은 만기 이전에는 이자만 상환하다가 만기에 일시로 원금을 상환하는 방식이다.
  \item 체증분할상환방식은 원리금 상환액 부담을 초기에는 적게 하는 대신 시간이 경과할수록 원리금 상환액 부담을 늘려가는 상환방식이다.
  \item 원리금균등분할상환방식은 원금이 상환됨에 따라 매기 이자액의 비중은 점차적으로 줄고 매기 원금상환액 비중은 점차적으로 증가한다.
\end{choices}
\end{question}
\answer{2}
